% ==== CODE INPUT: new, needs review ====
% Transferred from the confirmation report (paper.tex, Chapter 4 "Exploring techniques for learning with noisy labels for QSAR", and the opening of Chapter 5), shortened, with the "Change:" bullets applied.
This section reuses the exploratory study of learning with noisy labels (LNL) first reported in \citet{Punt2026confirmation}. In Chapter~\ref{chap:robustness}, the choice of representation explained more of the variance in clean performance than the choice of model. The relationship flipped for performance lost under noise: there, the model had the stronger effect. We therefore asked whether training methods built for noisy labels improve noise robustness, and whether the representation still matters under them.

LNL methods, and the five groups of \citet{Song2022}, are described in Section~\ref{sec:lnl_background}. They have been studied across machine learning, including QSAR \citep{Cortes2015}. \citet{Cortes2015} examined 12 algorithms on 12 bioactivity datasets. They added Gaussian noise ranging from zero up to each dataset's range of pIC50 values. Gradient boosting tolerated noise worst. SVM and GP models with polynomial kernels were sensitive at high noise in some cases. The other algorithms degraded smoothly and linearly with the noise level. Most LNL methods were built for image classification with flipped labels \citep{Song2022}. The labels here are continuous, so these methods do not transfer automatically. Cheminformatics data adds further challenges: high-dimensional molecular representations and a wide variety of chemical domains.

The confirmation report tested 12 LNL methods on the HOMO-LUMO gap from QM9, using 10,000 molecules and scaffold splits. Gaussian noise with standard deviation $\sigma$ from 0 to 1.0, in steps of 0.1, was added to the training labels only. Test labels stayed clean. The methods were paired with five representations: ECFP4 (2,048 bits), Sort \& Slice \citep{Dablander2024}, PDV, canonical SMILES and randomised SMILES. Depending on the method, it could take up to four optional add-ons. These were the model's own uncertainty, a distance between molecules, an alternative GP kernel, and an alternative loss function. The distances were Tanimoto on ECFP4, Euclidean on PDV, and Mahalanobis. The GP kernels were Tanimoto, RBF and Mat\'{e}rn. In total we tested 221 configurations over method, loss function, distance and representation. The report grouped the 12 methods into five categories of its own: meta-learning, sample selection, curriculum learning, iterative refinement and GP extensions. These are not the five groups of \citet{Song2022}. Four of the 12 remove suspected molecules and then retrain: small-loss selection, confident learning, distance-based selection and multistage cleaning. These four are noise mitigation. The other eight are noise-robust training.

We measured noise robustness by retention: $R^2$ under noise divided by $R^2$ with clean labels, as a percentage. Methods that used a distance did worse than those that did not. All 20 configurations with the lowest retention used a distance. None of the configurations with a distance reached both a clean $R^2$ above 0.6 and a retention above 80\%. The failure sits in the distance step. The same methods without the distance add-on did not fail this way. Not all use of similarity failed, however. The heteroscedastic GP compares molecules through its kernel and still held up. The failure is specific to judging one label by its neighbours' labels.

The methods that held up estimate noise from the model's own loss or uncertainty. DivideMix \citep{li2020dividemix} kept the most, with 86.7\% retention (standard deviation 0.4\%). Small-loss selection \citep{Han2018} kept 84.2\% (standard deviation 2.1\%). The uncertainty curriculum \citep{Bengio2009} kept 82.7\% (standard deviation 3.5\%). The heteroscedastic GP \citep{liu2020} kept 79.0\%, averaged over 15 configurations (3 kernels and 5 representations). Its $R^2$ was 0.71 with clean labels and 0.56 at $\sigma = 0.6$, the noise level its retention was measured at. The report does not state the noise level for the other methods' retention. Regularisation alone kept less: 73.0\% for Mixup and 74.0\% for sharpness-aware minimisation. None of the methods that held up compares a label with its neighbours' labels. Under these methods, the choice of representation explained even less variance than in Chapter~\ref{chap:robustness}. Chapter~\ref{chap:similarity} compares twelve LNL methods on the current datasets.

The study used one dataset, one property and one noise type: Gaussian noise on every training label. Chapter~\ref{chap:similarity} repeats the question on more datasets and noise types. The distance step was also added to methods designed without one. A poor adaptation is therefore an alternative explanation. The recommendation of noise-robust training over noise mitigation (Chapter~\ref{chap:discussion}, Limitations) starts here.

There are two possible reasons for the poor performance of the distances. First, a single molecular representation may lack sufficient information for distance calculations, with and possibly even without noise. PDVs encode global molecular properties, ECFPs contain substructure information, and MHG-GNNs encode the molecular graph \citep{kishimoto2023}. These single sources of data may not be sufficient for robust similarity estimation. Second, some of these representations are sparse. They contain many features that are not relevant to the molecules of a given dataset. When a large proportion of features is uninformative, similarity measures become dominated by them. Each uninformative feature that varies adds one more term to the distance, so their share of it grows linearly with their number. The evidence here does not separate the two reasons. The second is the one we sought to address, with a representation of informative features and low dimensionality. The budget-constrained hybrid representation concatenates the top-ranked features by importance from a selection of base representations.
% ==== END CODE INPUT ====
